\documentclass[UTF8,fontset=none,11pt,a4paper]{ctexart}
\usepackage[left=21mm,right=21mm,top=27mm,bottom=22mm,headheight=26pt,headsep=6mm]{geometry}
\usepackage{fontspec}
\setmainfont{texgyretermes-regular.otf}[BoldFont=texgyretermes-bold.otf,ItalicFont=texgyretermes-italic.otf]
\setCJKmainfont{Songti.ttc}[Path=/System/Library/Fonts/Supplemental/,FontIndex=6,AutoFakeBold=2]
\setCJKsansfont{Songti.ttc}[Path=/System/Library/Fonts/Supplemental/,FontIndex=6,AutoFakeBold=2]
\usepackage[table]{xcolor}
\definecolor{TableHead}{HTML}{EEF0ED}
\definecolor{Muted}{HTML}{666C68}
\usepackage{booktabs,longtable,array,fancyhdr,needspace,xurl}
\usepackage[unicode,hidelinks]{hyperref}
\urlstyle{same}
\linespread{1.32}
\setlength{\parindent}{2em}
\setlength{\parskip}{3pt}
\setlength{\emergencystretch}{3em}
\setlength{\LTpre}{8pt}\setlength{\LTpost}{10pt}
\ctexset{section={format=\fontsize{15}{20}\selectfont\bfseries,beforeskip=18pt,afterskip=8pt},subsection={format=\fontsize{12}{16}\selectfont\bfseries,beforeskip=12pt,afterskip=6pt},subsubsection={format=\normalsize\bfseries,beforeskip=9pt,afterskip=4pt}}
\pagestyle{fancy}\fancyhf{}
\renewcommand{\headrulewidth}{.3pt}
\renewcommand{\footrulewidth}{0pt}
\fancyhead[L]{\parbox[b]{\linewidth}{\fontsize{8}{10}\selectfont\color{Muted}现金账户：银行存款还是政府货币基金·银行存款归集项目披露}}
\fancyfoot[C]{\small\color{Muted}\thepage}
\hypersetup{pdftitle={银行存款归集项目披露},pdfsubject={仅供模拟使用·现金账户：银行存款还是政府货币基金·20-01},pdfauthor={金融决策模拟材料}}
\begin{document}
{\noindent\fontsize{20}{25}\selectfont\bfseries 银行存款归集项目披露\par}
\vspace{7pt}{\noindent\small\color{Muted}仅供模拟使用·现金账户：银行存款还是政府货币基金·20-01\par}\vspace{10pt}
\section*{定义}
本披露中的“平台”“我们”按语境包括经纪服务公司、清算服务公司及关联公司。“您”或“账户所有人”指开户申请列明的所有人。存入账户、收到而待投资、待提款或待进一步指示的款项称为“\allowbreak\mbox{现金余额”。}


项目银行存款是现金余额已经归集至参与银行形成的存款；货币市场基金溢出部分不属于项目银行存款。适用于不同现金管理、退休及健康储蓄账户的银行清单\allowbreak\mbox{可能不同。}


\section*{概述}
选择存款归集项目后，账户现金余额自动存入或从一个或多个参与银行的计息存款账户取出。参与银行不可接受存款，或现金余额超过项目可用容量及限额时，有关现金归集至货币市场基金\allowbreak\mbox{溢出安排。}


是否选择参加项目，由开户申请或现有账户变更指示决定。平台不替您决定是否采用该核心头寸，也不替您接受银行清单及其变更。您可以随时将现金余额移出项目，或停止采用项目，而不承担项目\allowbreak\mbox{退出罚金。}


无法存入银行，或平台判断可能不能适用存款保险的余额，按下文规则放入货币基金溢出安排。这会改变有关余额的保险适用、收益计算和\allowbreak\mbox{存取机制。}


\textbf{归集安排不作为长期投资选择。}账户有其他可用核心头寸，具体可用性按账户\allowbreak\mbox{协议处理。}


\section*{项目运作——存款}
每个营业日，账户现金余额自动归集至平台代表您及其他参加客户在参与银行设立的项目存款账户。营业日不包括银行假日及纽约证券交易所休市日，例如耶\allowbreak\mbox{稣受难日。}


归集从账户清单中的主要参与银行开始，直至该行的项目存款达到最高存款限额。剩余现金依清单顺序进入下一个可接受存款的银行，直至没有剩余现金，或各可用银行均已\allowbreak\mbox{达到限额。}


各可用银行均达到限额后，剩余现金进入货币市场基金溢出安排。在个别有限情形下，主要银行仍可能持有超过最高存款限额的余额，超出保险限额的部分不受美国联邦存款保险公司\allowbreak\mbox{（FDIC）保障。}


\subsection*{最高存款限额}
项目最高存款限额始终为当时适用的个人退休账户标准最高存款保险金额的98\%。例如标准保险金额为250,000美元时，项目最高存款限额为245,000美元。


这一项目配置限额与FDIC保险的实际合并计算是不同事项；本人的其他同银行同所有权类别存款\allowbreak\mbox{仍须合并。}


\subsection*{提款}
营业日结束时，账户需要资金支付借项，先使用货币基金溢出余额。仍需要资金，或没有溢出余额时，项目银行存款自动在下一营业日归\allowbreak\mbox{集回账户。}


某些日间实际或预计交易产生的借项，改为使用当日赎回溢出基金或当日提取银行存款的\allowbreak\mbox{款项结算。}


\textbf{银行提款顺序与存入顺序相同：}先从主要银行提取，再顺次从后续可用银行提取，直至借项付清或全部项目银行存款已经取回。留在各银行的存款不因该过程重新分配。因此，主要银行可以低于配置限额，而其他银行仍处于\allowbreak\mbox{配置限额。}


\subsection*{利息入账}
银行收到项目存款当天开始计息，提取当天停止计息。账户报表分别列明每个参与银行的\allowbreak\mbox{应计利息。}


银行存款每日计息，每月复利，于每月最后一个营业日记入项目存款账户。利息入账后归集至您的账户，可形成新的现金余额，再按存款\allowbreak\mbox{规则归集。}


溢出基金通常每日宣布股息、每月支付股息。形成的现金余额同样按有关归集\allowbreak\mbox{规则处理。}


\subsection*{存款保险限额变更}
适用于个人退休账户或普通单一所有权存款账户的标准保险金额变更时，平台按变更生效日重新确定有关账户的项目\allowbreak\mbox{配置限额。}


限额提高时，平台可以提高一个或多个银行的余额，减少清单中需要使用的银行数；限额降低时，可将余额分配至后续银行。能够参与的银行容量及溢出机制\allowbreak\mbox{继续适用。}


\subsection*{所有权记录}
平台保存反映各客户权益的记录。银行账户可以以平台或清算代理的名义设立，但记录须反映客户的受益所有权。FDIC保险依客户的所有权类别及适用规则计算，而不是对平台名称下所有客户合计只给一份\allowbreak\mbox{保险限额。}


\section*{参与银行}
各参与银行可以设定其可接受的项目存款总额，即银行存款容量。达到容量后，该银行可以不再接受新的项目存款。在容量尚未用满时，银行也可能因为其他原因无法\allowbreak\mbox{接受存款。}


\subsection*{银行状态}
\textbf{可用：}原则上可以接受项目存款，除非适用以下\allowbreak\mbox{其他状态。}


\textbf{不可用：}不能接受额外存款。该状态本身不改变已经持有的存款，但阻止新增存款。在没有可用溢出安排，且该银行为清单中唯一银行，或主要银行之外各银行均不可用、已排除或待激活等特定情形下，文件允许\allowbreak\mbox{例外处理。}


\textbf{溢出银行：}暂不能担任主要银行，已经存入的存款不因此变化。“溢出银行”这一状态与将余额投入溢出货币市场基金的安排不同。该银行是有关主清单唯一银行时，可以\allowbreak\mbox{适用例外。}


\textbf{待激活：}已经加入主清单，但尚不能接受\allowbreak\mbox{项目存款。}


\textbf{客户已排除：}您决定不采用该银行。银行仍列在您的清单中，但不接受本账户新归集存款；这只影响您的账户，不改变该银行在整个项目中的参与状态。排除须联系\allowbreak\mbox{平台代表。}


\subsection*{账户银行清单}
平台每天根据有关账户类型的银行主清单生成账户银行清单。不同开户日期，或不同账户类型，可能对应\allowbreak\mbox{不同清单。}


清单中的次序决定银行使用次序。第一家为主要银行。平台保留提前通知后修改清单的权利；分配初始清单后继续使用账户，或收到变更通知后继续使用，表示\allowbreak\mbox{接受清单。}


账户建立后，可以按适用协议排除一家或多家银行；存在其他清单时，也可以申请采用\allowbreak\mbox{另一清单。}


\subsection*{新账户}
新开户人可以按开户申请及客户协议选择归集。平台可因任何原因暂停或恢复向新账户提供该项目；暂停时可以选择其他可用\allowbreak\mbox{核心头寸。}


同一天新建立的同类账户采用当日生成的初始清单。清单由有关主清单中的部分或全部银行组成，并依主清单变更\allowbreak\mbox{自动更新。}


\subsection*{现有账户转入归集项目}
符合资格的现有账户可以按相关客户协议，在可用货币基金核心头寸与银行归集之间变更。决定转入归集时，采用当日分配给新账户的清单，并可依文件规定\allowbreak\mbox{排除银行。}


\subsection*{已采用项目的账户}
登入账户的持仓页面，可以查看当前清单、每行项目存款及银行状态。有银行状态、清单或项目相关问题时，可联系\allowbreak\mbox{平台代表。}


\subsection*{客户发起清单变更}
参加项目至少须保留一家可用银行。排除某银行使其在本账户中不再接受存款，但仍显示其状态。该排除指示可以撤回，只要银行仍参加项目且可以\allowbreak\mbox{接受存款。}


如果申请换清单，原已排除银行的指示原则上延续。但不能排除新清单的主要银行；要维持该排除，须使用主要银行不同的\allowbreak\mbox{另一清单。}


排除仍有余额的银行时，该余额被视为账户现金，并按规则重新分配至其他银行或溢出基金。改用另一清单时，原全部项目存款也按新清单\allowbreak\mbox{重新分配。}


\subsection*{银行加入及退出}
新银行可以加入清单，原银行可以退出。单纯增加或删除银行不重新排列其余银行的顺序。退出后，该银行不再接受新存款；平台提供更新清单，并将原存款转到其余银行或\allowbreak\mbox{溢出安排。}


主要银行退出时，从清单中指定另一可用银行为主要银行。每个参与账户至少须有一家可用银行；银行退出与客户排除指示共同导致无可用银行时，现金进入\allowbreak\mbox{溢出基金。}


\subsection*{银行状态变更}
新增银行可能先显示待激活，通常在加入清单约三十日后才能接受存款；客户可在待激活期间排\allowbreak\mbox{除该银行。}


银行由待激活、不可用或客户已排除状态恢复可用后，主要银行超过最高配置限额的余额可自动转往恢复可用的银行，按存款\allowbreak\mbox{次序处理。}


\subsection*{其他项目变更}
条件允许时，平台提前通知影响账户的变更。银行不愿或不能接受存款、项目参与终止，或银行持续经营及保险资格出现疑问等特殊情况下，您指示平台可以不事先通知，在日间修改清单、限制归集金额，或将余额移至其他银行或\allowbreak\mbox{溢出基金。}


没有可用银行时，部分或全部现金可以进入溢出基金。收到变更通知后继续使用账户，表示\allowbreak\mbox{同意变更。}


平台也可能必须把核心头寸改为另一可用核心头寸。提前书面通知后，您未在规定期间内提出不同指示，视为同意并指示平台提取银行及溢出余额、转入新\allowbreak\mbox{核心头寸。}


新核心头寸可能不适用FDIC保险，而可能适用美国证券投资者保护公司（SIPC）保护。平台不保证新头寸提供相同或更高收益；实际收益可能较低，平台所得报酬也可能不同\allowbreak\mbox{或更高。}


签署开户申请或指示采用归集，表示已经收到及阅读本披露，并同意必要时变更核心头寸。转入独立第三方顾问或平台关联顾问的有关服务时，可以按文件终止参与并清算\allowbreak\mbox{有关余额。}


您同意平台向项目服务方提供姓名、社会保障号码或税务号码、地址及出生日期等资料。银行等服务方仅为履行其本身或平台的法定及监管义务而使用这\allowbreak\mbox{些资料。}


\section*{现金余额的访问}
只能通过本平台账户访问现金，不能直接联系参与银行提取您的项目存款。送达平台的征税、扣押等法律程序可以作用于\allowbreak\mbox{有关存款。}


银行存款是银行对您的直接债务，不是平台的债务。平台不保证参与银行的财务状况。联邦银行规则允许银行在提款前要求七日通知。您在项目存款账户的权益\allowbreak\mbox{不能转让。}


\subsection*{存款账户}
参与银行采用活期存款账户持有项目存款。您在银行没有独立活期子账户，存款与其他项目客户存款合并持有。文件为便于描述而将分配给您的部分称为“您的”活\allowbreak\mbox{期子账户。}


\subsection*{报表及确认}
每期报表显示核心头寸余额、各银行在报表期最后营业日的存款及溢出基金余额，列明归集进出交易，并分别列明银行及溢出收益。报表替代每笔归集的\allowbreak\mbox{单独确认。}


您负责监测每个银行全部同类别存款，包括通过其他安排存入的款项。银行清单变化可能影响实际\allowbreak\mbox{保险覆盖。}


\subsection*{货币基金溢出}
各可用银行达到配置限额，或项目及银行没有足够存款容量时，不能存入的现金进入\allowbreak\mbox{溢出基金。}


账户借项先使用溢出基金，再从主要银行及后续银行提款。某些日间实际或预计交易使用当日赎回及提款款项；其他借项按下一营业日提款\allowbreak\mbox{安排结算。}


银行恢复容量之后，已在溢出基金中的款项不会因此自动转回或重新平衡到银行。除用于支付借项、提款或适用再平衡事件外，仍留在\allowbreak\mbox{溢出基金。}


持仓页面显示溢出余额；报表显示相关余额、归集记录及收益。溢出部分不适用FDIC存款保险，可以在适用限额内受SIPC保护。适用时提供有关溢出基\allowbreak\mbox{金说明书。}


\section*{银行存款利率}
利率可以随时改变，并受到银行支付利率、平台费用、账户类型及文件列明的其他因素影响。项目利率可能低于直接在参与银行开户的利率，或其他非FDIC投资及货币\allowbreak\mbox{基金收益。}


不同账户类型可采用不同利率；本人其他账户取得的利率不必与本账户一致。具体当时利率可向平台或利率\allowbreak\mbox{页面查询。}


\section*{溢出基金收益}
溢出基金通常采用七日年化收益指标，随基金投资收入及其他因素变化。有效收益率还考虑\allowbreak\mbox{年度复利。}


货币基金取得利息、股息及其他收入，扣除费用后向持有人分配；实现资本收益扣除损失后，也可以分配。通常每日宣布、每月支付股息。收到资金日开始累计股息，提\allowbreak\mbox{取日停止。}


\section*{FDIC保险及SIPC保护}
项目存款须与本人在同一银行、同一所有权类别中的其他存款合并计算，包括其他归集项目、储蓄、支票、货币市场存款及银行直接\allowbreak\mbox{发行存单。}


单一所有权账户通常最高250,000美元；共同账户每位共同所有人份额通常最高250,000美元。退休账户和健康储蓄账户依各自所有权分类处理；信托存款有另外规则。不同所有权类别依法可以\allowbreak\mbox{分别保障。}


平台不掌握本人全部非项目存款，也不负责替本人监测是否超过保险限额。本人须监测合并金额。超过适用限额的资金不受FDIC保险。


保险金额包括符合规则的本金及应计未付利息。银行倒闭需要保险支付时，由FDIC提供；没有特定必须向本人提供款项的时间期限，资金访问可能显著延迟，也可能要求本人向FDIC及平台\allowbreak\mbox{提交文件。}


账户证券及溢出基金是投资产品，不适用FDIC保险，没有银行或政府本金保证，可能损失本金。SIPC对客户资产保护的适用上限为500,000美元，其中现金请求上限250,000美元；不赔偿投资市场\allowbreak\mbox{价值波动。}


现金只有在成为参与银行持有的项目存款之后才适用FDIC。在平台持有或银行进出途中，一般按SIPC规则处理。另选货币基金或其他核心投资头寸时，也不因此成为FDIC存款。


\section*{平台与其他机构的关系}
平台为经纪账户提供相关服务；清算服务公司作为代理在银行设立合并项目账户，并办理归集。项目管理方提供记录及技术服务。退休或健康储蓄账户如适用，另有\allowbreak\mbox{托管主体。}


平台可以同时是项目管理方或银行的客户，并存在其他金融往来；项目管理方也可以是参与银行。以后关联公司可以被指定为项目管理方、参与银\allowbreak\mbox{行或两者。}


平台母公司对某参与银行持有少数、非控制性权益。关联基金管理公司负责管理溢出基金，并取得管理及\allowbreak\mbox{其他费用。}


\subsection*{平台与其他机构取得的利益}
银行向平台支付项目费用，通常按每日项目存款合计余额计算。费用因银行、账户类型及时间而异，存款增加通常使费用增加；同一银行在其他类似项目支付的费用也\allowbreak\mbox{可能不同。}


费用可达到该银行有关项目存款余额的年化4\%。平台可降低或免除部分费用；超过4\%的提高会另行通知。该费用由参与银行支\allowbreak\mbox{付给平台。}


银行付给平台的费用，按银行合同“费率”乘以存款余额，再减银行向项目账户支付的全部利息计算。合同费率通常与有效联邦基金利率或目标利率有关，可能按余额分档，并处于相关联邦基金利率上下0.75个百分点\allowbreak\mbox{的范围内。}


项目管理方为技术、会计及监管协助服务收取平台支付的年费，按每日项目存款合计余额的0.010\%计算，并设固定上限。管理方同时是参与银行时，银行付给平台的费用与平台付给管理方的费用\allowbreak\mbox{各自独立。}


银行将项目存款用于现有及新增贷款和投资。其收入来自这些资产收入与支付给存款人利息及平台费用之间的差额。银行可能支付低于市场或直接开户存款的利率。平台还可能因关联持股取得股\allowbreak\mbox{息或利润。}


溢出基金管理机构取得基金管理和运营费用，部分费用可以与\allowbreak\mbox{平台分享。}

\end{document}
